\section{High-rate qLDPC codes}
\label{app:high-rate-qLDPC-codes}

In this section, we give the definitions of the qLDPC codes used in this work.

\paragraph{Bivariate Bicycle (BB) codes}
The construction of BB codes~\cite{bravyi_high-threshold_2024,yoder_tour_2025} begins with the $\ell\times\ell$ cyclic shift matrix $S_\ell$, where the $i$-th row contains a single nonzero entry in column $(i+1) \bmod \ell$.
For example, $S_3$ is given by:
\begin{equation}
    S_3 = \begin{pmatrix}
        0 & 1 & 0 \\
        0 & 0 & 1 \\
        1 & 0 & 0
    \end{pmatrix}.
\end{equation}
We define two generator matrices $x = S_\ell \otimes I_m$ and $y = I_\ell \otimes S_m$, where $\otimes$ denotes the Kronecker product.
These generators satisfy $x^{\ell} = y^{m} = I_{\ell m}$, and any monomial $x^p y^q$ (for $0 \leq p < \ell$ and $0 \leq q < m$) corresponds to the matrix $S_\ell^p \otimes S_m^q$, allowing us to treat monomials and matrices interchangeably.
We define the transpose of $x^p y^q$ to be $(x^p y^q)^\top = x^{\ell -p} y^{m - q} = x^{-p} y^{-q}$.

A BB code is parameterized by two polynomials $A$ and $B$, each expressed as a sum of three distinct monomials:
\begin{equation}
    A = A_1 + A_2 + A_3, \quad B = B_1 + B_2 + B_3,
\end{equation}
where each $A_i$ and $B_i$ is a distinct monomial of the form $x^p y^q$.
The parity-check matrices for the $X$ and $Z$ stabilizers are then constructed as:
\begin{equation}
H_X = [A|B], \quad H_Z = [B^\top|A^\top],
\end{equation}
where $H_X$ and $H_Z$ are $n/2 \times n$ matrices with $n = 2 \ell m$ qubits.

A widely known example of BB codes is the $[[144,12,12]]$ code,
which is defined by $ \ell = 12 $, $ m = 6 $, and the polynomials:
\begin{equation}
    A=1+y+x^{3} y^{-1}, \quad B=1+x+x^{-1} y^{-3}.
\end{equation}
% We use the logical representatives following Ref.~\cite[Eqs. (31)--(34)]{yoder_tour_2025} and syndrome extraction circuit in Ref.~\cite{bravyi_high-threshold_2024}.
We use the logical Pauli representatives~\cite[Sec.~A.1]{yoder_tour_2025}
that has undergone a series of optimizations~\cite{bravyi_high-threshold_2024,
cross_improved_2024,williamson_low-overhead_2024,yoder_tour_2025}
targeting logical operations with low-overhead ancilla system.
Fortunately, such representatives are also friendly to injection,
where overlaps are intentionally designed to be limited,
and the weight of the logical Pauli representatives is exactly $ d=12 $.

\paragraph{Hypergraph Product (HGP) Codes}

We also inject on the $[[225,9,4]]$ HGP code 
from~\cite{xu_constant-overhead_2023} to demonstrate the generality of our scheme.
Such family of HGP codes are constructed as the product of classical LDPC codes defined by $(3,4)$-regular Tanner graphs, and have good circuit-level performance with constant overhead~\cite{xu_constant-overhead_2023},
and fast, parallelizable logical computation~\cite{xu_fast_2024}.

Denote the parity check matrices of the two codes as $H_1$ and $H_2$, with shape $ r_i \times n_i $ respectively.
The parity check matrix of the HGP code is given by:

\begin{equation}
    \begin{aligned}
    & H_X=\left(\begin{array}{ll}
    H_1 \otimes I_{n_2} & I_{r_1} \otimes H_2^\top
    \end{array}\right) \\
    & H_Z=\left(\begin{array}{ll}
    I_{n_1} \otimes H_2 & H_1^\top \otimes I_{r_2}
    \end{array}\right),
    \end{aligned}
\end{equation}
where $ I_{m} $ is the $ m $-dimensional identity matrix.

The standard logical Pauli representatives construction for QEC codes~\cite[Sec.~4.1]{gottesman_stabilizer_1997}
don't guarantee short-weight logical Pauli representatives for all the logical qubits.
For HGP codes, we construct the logical Pauli representatives with shorter weight than the generic construction.
We use the syndrome extraction circuit in Ref.~\cite{xu_constant-overhead_2023}.

\section{Maximal injectable sets}
\label{sec:maximal-injectable-sets}

An index set $\cI \subseteq [k]$ is injectable under injection sites
$\bm Q=\{Q_1,\dots,Q_k\}$ with $Q_i\in O_i$ if it satisfies
injection site independence (\Cref{def:injection-site-independence})
and compatible pairing (\Cref{def:generic-compatible}).

We compute injectable sets by enumerating injection-site choices and then checking subset feasibility:
\begin{enumerate}
\item Enumerate the Cartesian product $\prod_i O_i$ to obtain candidate $\bm Q$.

\item For fixed $\bm Q$, build a graph on vertices $\{1,\dots,k\}$, where edge $(i,j)$ means $i,j$ cannot be injected together because injection site independence is violated.
Any injectable set must be an independent set in this graph.
Compute all maximal independent sets (MISs) as candidates.

\item For each MIS candidate, enumerate subsets and test compatible pairing using signatures.
For each physical qubit $q$ in any $C_{i,j}$ (defined in \Cref{eq:unified-Cij}), define its signature
\begin{equation}
\sigma(q) \;:=\; \bigl(\mathbf 1[q\in C_{i,j}]\bigr)_{i,j\in\cI}.
\label{eq:signature-def}
\end{equation}
Compatible pairing exists iff every signature class $\{q:\sigma(q)=s\}$ has even cardinality.
Within each signature class, any disjoint pairing is valid.
\end{enumerate}

\section{Noise model}

The noise model we used in this work is the same as the noise model in~\cite{li_magic_2015,lao_magic_2022}, as shown in \Cref{tab:li_noise_channels,tab:li_noise_model}.

\begin{table}[t]
    \centering
    % \renewcommand{\arraystretch}{1.25}
    \begin{tabular}{c|l}
    \toprule
    \textbf{Noise channel} & \textbf{Definition} \\
    \midrule
    $\mathrm{MERR}(p)$ &
    $m \mapsto m\oplus e,\ \ e\sim \mathrm{Bernoulli}(p)$ \\
    % \hline
    $\mathrm{XERR}(p)$ &
    $\rho \mapsto (1-p)\rho + p\,X\rho X$ \\
    % \hline
    $\mathrm{ZERR}(p)$ &
    $\rho \mapsto (1-p)\rho + p\,Z\rho Z$ \\
    % \hline
    $\mathrm{DEP1}(p)$ &
    $\rho \mapsto (1-p)\rho + \frac{p}{3}\sum_{P\in\{X,Y,Z\}} P\rho P$ \\
    % \hline
    $\mathrm{DEP2}(p)$ &
    $\rho \mapsto (1-p)\rho + \frac{p}{15}\!\!\sum_{P\in\{I,X,Y,Z\}^{\otimes2}\setminus\{II\}}\!\! P\rho P$ \\
    % \hline
    \bottomrule
\end{tabular}
\caption{Definitions of the noise channels.}
\label{tab:li_noise_channels}
\end{table}

 
\begin{table}[t]
\centering
\renewcommand{\arraystretch}{1.2}
\begin{tabular}{l|l}
\toprule
\textbf{Ideal operation} & \textbf{Noisy operation} \\
\midrule
% \hline
$R_X$ (prepare ${+}1$ eigenstate of $X$)
& $R_X \circ \mathrm{ZERR}(p_{\rm IN})$ \\
$R_Y$ (prepare ${+}1$ eigenstate of $Y$)
& $R_Y \circ \mathrm{XERR}(p_{\rm IN})$ \\
$R_Z$ (prepare ${+}1$ eigenstate of $Z$)
& $R_Z \circ \mathrm{XERR}(p_{\rm IN})$ \\
% \hline
$M_Z$ (measure $Z$)
& $M_Z \circ \mathrm{MERR}(p_{\rm M})$ \\
% \hline
$M_X$ (measure $X$)
& $M_X \circ \mathrm{MERR}(p_{\rm M})$ \\
% \hline
$U_1$ (any single-qubit unitary)
& $U_1 \circ \mathrm{DEP1}(p_{\rm 1Q})$ \\
% \hline
$U_2$ (any two-qubit unitary)
& $U_2 \circ \mathrm{DEP2}(p_{\rm 2Q})$ \\
$\text{IDLE}$ (idle)
& $\text{IDLE} \circ \mathrm{DEP1}(p_{\rm IDLE})$ \\
% \hline
\bottomrule
\end{tabular}
\caption{The noise model we used in this work, the same as circuit-level noise model in~\cite{li_magic_2015,lao_magic_2022}. The notation $G \circ \mathcal{E}$ indicates that the ideal operation $G$ is \emph{followed} by the noise channel $\mathcal{E}$.}
\label{tab:li_noise_model}
\end{table}
